\documentclass[11pt]{article}
\usepackage[margin=25mm]{geometry}
\usepackage{fontspec,unicode-math}
\setmainfont{Latin Modern Roman}
\setmathfont{Latin Modern Math}
\usepackage{amsmath,amssymb,mathtools}
\usepackage{xurl,hyperref}
\hypersetup{colorlinks=true,urlcolor=blue,linkcolor=blue}
\setcounter{secnumdepth}{0}
\setlength{\parindent}{0pt}
\setlength{\parskip}{5pt}
\setlength{\emergencystretch}{3em}
\newcommand{\sref}[2]{\hyperlink{source:#1:#2}{[S#2]}}
\newcommand{\cataloguescope}[1]{\paragraph{Source question.}#1}
\newcommand{\floor}[1]{\left\lfloor#1\right\rfloor}
\begin{document}
\section{Approximating Fare-Zone Assignment on Trees}\label{statement}

\subsection{Definitions and mathematical statement}

\paragraph{Shared notation.}Standard number systems are denoted by $\mathbb N,\mathbb Z,\mathbb Q,\mathbb R,\mathbb C$; index ranges, computational models and specialized notation follow the cited problem source and the definitions below.


$y_e=1$ selects an edge as a fare-zone boundary. For a commodity path $P_i$, the constraint limits how many selected edges lie along it; the objective adds its weight once for each selected edge it traverses.
\paragraph{Statement recorded in the supplied source.}
Given a tree $G=(V,E)$ and weighted commodity paths $(P_i,w_i,u_i)$, choose $y_e\in\{0,1\}$ to maximize $\sum_{e\in E}y_e\sum_{i:e\in P_i}w_i$ subject to $\sum_{e\in P_i}y_e\leq u_i$ for every $i$. What approximation guarantees are achievable on trees, particularly when the maximum number of paths sharing an edge is bounded by a parameter $k$? \sref{statement}{2}
\paragraph{Scope qualification.}
The input includes the listed path weights and budgets; no omitted transit-network interpretation changes the displayed optimization problem.
\paragraph{Source status.}
The supplied catalogue labels this record partially solved. Its source-status metadata was checked against the supplied snapshot on 2026-10-11. This import does not independently establish current open status or mathematical correctness. The catalogue literature assessment is dated 2026-08-22.
\subsection{Short English statement}
Approximate the stated fare-zone optimization on trees, especially when few commodity paths share any edge.
\subsection{Sources}
\begin{itemize}
\item[\hypertarget{source:statement:1}{S1}] Ulam AI and the UnsolvedMath Contributors, UnsolvedMath: A Curated Collection of Open Mathematics Problems, catalogue entry. Catalogue curation and metadata: CC BY 4.0. The supplied source question is retained; definitions and the English summary are editorial additions.
\url{https://huggingface.co/datasets/ulamai/UnsolvedMath}.
\item[\hypertarget{source:statement:2}{S2}] Combinatorial Optimization (2025), Oberwolfach Reports, DOI 10.4171/owr/2024/50, pp. 2949--3000
\url{https://doi.org/10.4171/owr/2024/50}.
\item[\hypertarget{source:statement:3}{S3}] OWR report
\url{https://oa.tib.eu/renate/bitstreams/706cca58-da6b-4d8c-a979-ebf6e92658f0/download}.
\item[\hypertarget{source:statement:4}{S4}] Hoefer--Kauther--Pabst--Peis--Tran
\url{https://arxiv.org/abs/2512.19493}.
\textit{Underlying Oberwolfach Reports material retains its source-specific licence, including CC BY-SA 4.0 where applicable; see THIRD\_PARTY\_NOTICES.md.}
\end{itemize}
\paragraph{Preparation and review.}
Original question attributed to the cited source and UnsolvedMath contributors. Codex (GPT-6; exact serving revision not exposed) assisted with definitions, English summaries, source inspection and formatting. No new human mathematical authorship or independent proof verification is asserted. Catalogue identity checked against the complete supplied registry and existing repository statements; known flagged repetitions and obvious identical conjectures excluded. Full mathematical-equivalence/literature review has not been performed. Statement preparation and final review followed contribution prompts C and D; automated checks verify structure and identity, not open status.
\label{end:statement}
\end{document}
